\ssr{ЗАКЛЮЧЕНИЕ}

В ходе работы были исследованы два алгоритма для решения задачи коммивояжёра:
полный перебор и муравьиный алгоритм. Каждый из них имеет свои преимущества и
ограничения, которые определяют их применимость в различных ситуациях.

Метод полного перебора обеспечивает нахождение оптимального решения,
поскольку он проверяет все возможные перестановки маршрутов.
Однако из-за факториальной сложности его использование ограничено задачами с
числом городов <= 6.

Муравьиный алгоритм представляет собой эвристический метод,
который вдохновлён поведением муравьёв в природе. Этот алгоритм обеспечивает
приемлемые результаты для графов с 7 и более вершин. Несмотря на то, что он не гарантирует
нахождение оптимального решения, получаемые маршруты близки к оптимальным. Эффективность
алгоритма также зависит от его параметров. Результат исследования показал, что оптимальным
набором является $\alpha$ = 0.5, $\rho$ = 0.9, $t\_{max}$ = 500.

В ходе выполненной работы были выполнены следующие задачи:
\begin{itemize}
    \item[---] описаны стандартный метод решения задачи коммивояжера и метод на базе муравьиного алгоритма;
	\item[---] описана модель вычислений;
    \item[---] выполнена оценка трудоемкости разрабатываемых алгоритмов;
    \item[---] реализованы алгоритмы;
    \item[---] проведен сравнительный анализ двух рассмотренных методов решения задачи коммивояжера;
    \item[---] выполнена параметризация метода, основанного на муравьином алгоритме по трём его параметрам.
\end{itemize}